\section{Mathematical Model}
\label{sec:math model}

Based on the physical considerations discussed above, the electromechanical behavior 
of the material is described at the continuum level by a coupled system of partial differential equations.

The mechanical response is governed by the equations of elasticity,

\begin{equation}
    \label{eq:conductivity}
    \nabla \cdot \boldsymbol{\sigma}(u) = 0,
\end{equation}

where $u$ denotes the displacement field and $\boldsymbol{\sigma}(u)$ is the stress tensor, 
which depends on the strain through a constitutive relation.

The electrical response is described by the conductivity equation,

\begin{equation}
    \label{eq:1d-cond}
    \nabla \cdot \left( \sigma(\varepsilon(u)) \nabla V \right) = 0,
\end{equation}

where $V$ is the electrical potential and the conductivity $\sigma$ depends on the strain tensor 
$\varepsilon(u)$.

Together, these equations define a coupled system in which the mechanical field $u$ determines 
the strain $\varepsilon(u)$, which in turn modifies the conductivity $\sigma$, and thereby 
influences the electrical potential $V$.

From a physical perspective, this coupling provides an effective continuum description of how 
mechanical deformation alters the underlying electronic transport properties of the material.

The system is nonlinear due to the dependence of the conductivity on the strain field.

This formulation serves as the basis for the computational modeling of the piezoresistive effect, 
linking continuum mechanics with strain-dependent electrical conduction.
